\documentclass[14pt, a4paper]{extarticle}

% Підтримка української мови та кирилиці
\usepackage[T2A]{fontenc}
\usepackage[utf8]{inputenc}
\usepackage[ukrainian]{babel}
\usepackage{amsmath, amssymb}

% Налаштування полів (робимо трохи ширші для зручності тримання в руках)
\usepackage[left=2.5cm, right=2cm, top=2cm, bottom=2cm]{geometry}

% Півторний інтервал для зручного читання з аркуша
\usepackage{setspace}
\onehalfspacing

% Абзацний відступ
\usepackage{indentfirst}
\setlength{\parindent}{1.25cm}

% Вимкнення нумерації сторінок (за бажанням)
\pagestyle{plain}

\begin{document}
	
	\begin{center}
		\textbf{\Large Текст доповіді до дипломної роботи} \\
		\vspace{0.5cm}
		\large Тема: «Дослідження та програмна реалізація методу масштабування зображень з використанням барицентричної формули на основі вузлів інтерполяції Чебишева»
	\end{center}
	
	\vspace{1cm}
	
	\textbf{[СЛАЙД 1-2] Формулювання теми та актуальність}
	
	Задачі масштабування зображень є критично важливими у сферах геоінформаційних систем (ГІС), обробки супутникових знімків, медичної візуалізації та комп'ютерного зору. Від вибору алгоритму масштабування безпосередньо залежить баланс між обчислювальною складністю та якістю зображення.
	
	Актуальність нашого дослідження зумовлена тим, що з розвитком цифрових технологій зростають і вимоги до якості графіки. Нам постійно доводиться адаптувати зображення під різні формати екранів: від маленьких дисплеїв смартфонів до великих 4K-моніторів. При цьому в таких сферах, як медична візуалізація, супутникові знімки чи комп'ютерний зір, втрата навіть дрібних деталей під час зміни розміру картинки є критичною.
	
	Проблема полягає в тому, що більшість існуючих стандартних алгоритмів (таких як білінійна чи бікубічна інтерполяція) є компромісними. Вони працюють швидко, але часто призводять до розмиття зображення або появи артефактів на різких контурах об'єктів. З іншого боку, складніші глобальні математичні методи раніше вважалися або занадто повільними для комп'ютерної графіки, або страждали від крайових спотворень, відомих як феномен Рунге. Вони не вирішують в повній мірі сучасні задачі.
	
	Наша робота полягає у спробі створення нового методу, який дозволить обійти ці недоліки, уникнути артефактів та зберегти високу структурну точність пікселів (особливо при сильному зменшенні зображень), це є надзвичайно актуальною задачею, яку ми і вирішуємо у цій роботі. Інтерполяційні формули за вузлами Чебишева широко використовуються в математиці, і також можуть бути використані для задачі масштабування зображень.
	
	Метод, який ми представляємо, працює як у напрямку зменшення, так і збільшення зображень, у тому числі для великих коефіцієнтів масштабування. Це неадаптивний метод, що належить до класу методів інтерполяції. Більшість класичних методів зазвичай базуються на рівномірних сітках вузлів і розглядають лише локальну інтерполяцію, а не глобальну.
	\vspace{0.5cm}
	\textbf{[СЛАЙД 3] Математична модель}
	
	З теоретичної точки зору, будь-яке зображення можна представити як двовимірну неперервну функцію просторових координат, задану на квадраті від мінус одиниці до одиниці. Відповідно, цифрове зображення — це просто дискретний набір значень цієї функції, обчислених у певних точках, тобто на сітці вузлів.
	
	У цій роботі для формування такої сітки ми на відміну від статті що була взята за основу використовуємо не нулі Чебишева а екстремуми поліномів Чебишева другого роду. Наша задача масштабування зводиться до пошуку апроксимованих значень функції на новій сітці: більш щільній, якщо ми збільшуємо зображення, або розрідженій, якщо зменшуємо.
	
	Головною математичною особливістю нашого підходу є те, що для максимально ефективної реалізації інтерполяції ми узагальнили відому одновимірну другу барицентричну формулу на двовимірний випадок. Таке узагальнення стало практично можливим завдяки тому, що для побудови барицентричної формули є аналітичні формули для розрахунку коефіцієнтів.
	Варто відзначити, що сама по собі барицентрична формула є універсальною і може бути використана для будь-яких інших систем вузлів. Однак, у такому разі побудова інтерполяційного полінома та розрахунки коефіцієнтів були б значно важчими та ресурсоємнішими. Використання ж саме запропонованих вузлів дозволяє зробити ці обчислення значно простіше.
	\vspace{0.5cm}
	\textbf{[СЛАЙД 4] Програмна реалізація}
	
	Для практичного підтвердження ефективності математичної моделі було розроблено спеціалізоване програмне забезпечення. Програму реалізовано мовою C++ з використанням графічного фреймворку Qt. Використання Qt дозволило створити зручний інтерфейс користувача для завантаження зображень, вибору коефіцієнтів масштабування та візуального порівняння результатів.
	
	Щоб об'єктивно оцінити наш алгоритм, ми порівняли його з еталонними методами з популярної бібліотеки OpenCV, такими як метод найближчого сусіда, білінійна та бікубічна інтерполяція, а також метод Ланцоша. Оскільки глобальна двовимірна інтерполяція є обчислювально важкою задачею, ми оптимізували програму за допомогою засобів багатопотоковості мови C++. Це дозволило розподілити обчислення між ядрами процесора і суттєво зменшити час обробки пікселів.
	
	Якість масштабування оцінювалася за допомогою стандартних метрик: PSNR (пікове відношення сигналу до шуму) та SSIM (індекс структурної подібності). Експерименти показали, що при збільшенні зображень наш метод демонструє якість, порівнянну з класичною бікубічною інтерполяцією. Водночас, при зменшенні зображень він забезпечує значно кращі показники PSNR та SSIM, мінімізуючи втрату деталей.
	
	\vspace{0.5cm}
	\textbf{[СЛАЙД 5-6] Висновки}
	
	Підсумовуючи, у нашій роботі було досліджено та практично реалізовано сучасний математичний підхід до проблеми зміни розміру зображення, що базується на глобальній інтерполяції. Сама теоретична модель була запропонована в статті що була взята за основу, проте нашою головною метою та досягненням стала оптимізація та адаптація нового обчислювально важкого алгоритму для практичного використання. Відповідно до цієї моделі, ми виконуємо апроксимацію зображення до бажаного розміру за допомогою процесів глобальної інтерполяції на основі сіток Чебишева. Зокрема, представлений метод бере значення пікселів будь-якого вхідного зображення для побудови відповідного двовимірного інтерполяційного полінома на основі другої барицентричної формули, обчислення якого дає на виході масштабоване зображення.
	
	Однією з головних переваг цього методу є його висока гнучкість під час роботи в обох напрямках масштабування. Запропонований математичний підхід залишає відкритим шлях для майбутніх розробок та можливих покращень.
	
	Подальшим розвитком програмного забезпечення буде оптимізація методу масштабування зображень з використанням другої барицентричної формули на основі вузлів інтерполяції Чебишева для отримання найкращих результатів по часу розрахунків, автоматизація розрахунку порівняльних характеристик, що призначені для чисельної оцінки якості масштабування та покращення інтерфейсу програмного забезпечення для більш комфортного та інтуїтиво зрозумілого користування.
	
\end{document}